\section{Legal Notices: Permissions \& Events}
\label{sec:permissions-events}

\subsection{Permissions: a compiled tree}
A module can declare a \texttt{permissions:} tree --- nested access
control permissions the module contributes. Each node has a \texttt{key}
(its slug) and a \texttt{name} (used to name the generated identifier,
e.g.\ \texttt{name: managePosts} $\to$ \texttt{ManagePostsPermission}).
The full identifier is either set explicitly via \texttt{fullKey:} or
computed as \texttt{\{parent.fullKey\}.\{key\}} during preprocessing; a
node that groups children and left its key to be derived gets a trailing
\texttt{.*} (e.g.\ \texttt{post.*}), signaling it covers itself and
everything below.

\begin{lstlisting}
permissions:
  - key: post
    name: managePosts
    children:
      - key: create
        name: createPost
      - key: publish
        name: publishPost
\end{lstlisting}

Compiling this (Go target) produces one exported var \textbf{per root
permission} --- not one var wrapping the whole forest, and no separate
flat constant per node either. Each var is a fully typed, nested struct:
every node embeds \texttt{emigo.Permission} (promoting its own
\texttt{Key}) and has one field per child, so an editor autocompletes
straight down it with zero risk of a typo'd key only failing at runtime:
\begin{lstlisting}
var ManagePostsPermission = struct {
  emigo.Permission
  CreatePost  emigo.Permission
  PublishPost emigo.Permission
}{ /* ... */ }

var ManagePostsPermissionList = []emigo.Permission{
  ManagePostsPermission.Permission,
  ManagePostsPermission.CreatePost,
  ManagePostsPermission.PublishPost,
}
\end{lstlisting}
\texttt{<Name>PermissionList} sits alongside every root, built by
referencing the root var's own fields rather than re-declaring separate
literals, so it can never drift out of sync with it --- one list per
root, not one combined list a caller has to filter back down. The JS/TS
compiler mirrors this shape as a \texttt{const ... as const} object (or
\texttt{Object.freeze(...)} without \texttt{--tags typescript}) instead of
a Go struct, with the same one-const-per-root, \texttt{<Name>PermissionList}
convention.

\subsection{Events: flat facts, typed payloads}
A module can declare an \texttt{events:} list --- flat, unlike
\texttt{permissions:}, since an event is a single fact (``a post got
published''), not a tree. Each event has a \texttt{key} (its wire name),
optional localized \texttt{name}/\texttt{description}, a \texttt{payload},
and \texttt{permissions} describing who's allowed to know it happened ---
a list of \texttt{with} blocks, OR'd together, each entry's \texttt{with}
list AND'd (\texttt{[\{with: [A,B]\}, \{with: [C,D]\}]} means
\texttt{(A and B) or (C and D)}).

\begin{lstlisting}
events:
  - key: postPublished
    permissions:
      - with: ["post.publish"]
      - with: ["post.*"]
    payload:
      fields:
        - {name: postId, type: string}
  - key: commentAdded
    payload:
      dto: CommentDto
  - key: postArchived
\end{lstlisting}

Three events, three payload shapes: \texttt{postPublished} declares
\texttt{fields} inline, compiled into a dedicated
\texttt{PostPublishedEventPayload} struct exactly like an action's
request body, complete with CLI flags; \texttt{commentAdded} points
\texttt{dto:} at an existing dto, so its payload type is a plain alias,
not a new struct; \texttt{postArchived} declares no payload at all, so
its payload type aliases straight to \texttt{interface\{\}}. One exported
var is generated per event (PascalCase, trailing \texttt{Event}), always
paired with a \texttt{<EventName>Payload} type and a
\texttt{New<EventName>(payload)} constructor, and every var is collected
into \texttt{AllEventsList} for anything that wants to walk the whole
catalog. The compiler rejects the module outright if two differently
spelled keys (\texttt{postPublished} vs.\ \texttt{post\_published})
would normalize to the same generated identifier, rather than letting one
silently shadow the other.
